\documentclass[11pt,a4paper]{amsart}
\usepackage[margin=2.6cm]{geometry}
\usepackage{amsmath,amssymb,amsthm,mathtools}
\usepackage{graphicx,booktabs}
\usepackage{xcolor}
\usepackage[colorlinks=true,linkcolor=blue!50!black,citecolor=blue!50!black,urlcolor=blue!50!black]{hyperref}
\newtheorem{theorem}{Theorem}
\newtheorem{lemma}[theorem]{Lemma}
\newtheorem{corollary}[theorem]{Corollary}
\newtheorem{proposition}[theorem]{Proposition}
\theoremstyle{definition}
\newtheorem{definition}[theorem]{Definition}
\newtheorem{conjecture}[theorem]{Conjecture}
\newtheorem{remark}[theorem]{Remark}
\newcommand{\R}{\mathbb{R}}
\newcommand{\Z}{\mathbb{Z}}
\newcommand{\Lat}{\Lambda}
\DeclareMathOperator{\covol}{covol}

\title[Three-body lattice energies in the limit of steep decay]{Which lattice minimises a steep three-body power-law energy?\\ A max--min triangle problem}
\author{DRAFT}
\date{\today}

\begin{document}
\begin{abstract}
For a lattice $\Lat\subset\R^d$ of unit covolume consider the three-body energy
$T_\nu(\Lat)=\sum_{x,y\in\Lat\setminus\{0\},\,x\neq y}\bigl(|x|\,|y|\,|x-y|\bigr)^{-\nu}$, i.e.\ the sum over all lattice triangles through the origin of the product of the side lengths to the power $-\nu$.
In $d=2$ this is, up to normalisation, the two-loop modular graph function $C_{s,s,s}(\tau)$ with $s=\nu/2$.
We prove that as $\nu\to\infty$ every family of minimisers of $T_\nu$ converges to the rectangular lattice with aspect ratio $\sqrt{(\sqrt{17}-1)/2}$; in particular neither the hexagonal nor the square lattice is optimal for large $\nu$.
The proof reduces the limit to a max--min problem for the smallest triangle product of a lattice, which we solve exactly.
In $d=3$ we prove, by a computer-assisted argument (branch and bound over Minkowski-reduced Gram matrices plus an interval-arithmetic local lemma), that the body-centred cubic lattice is the unique maximiser of the smallest triangle product, so that minimisers of $T_\nu$ converge to BCC.
For finite exponents we prove, again computer-assisted, that $C_{2,2,2}$, $C_{3,3,3}$ and $C_{4,4,4}$ attain their minimum only at $\tau=i$ (square lattice), whereas $C_{1,1,1}$ is minimal at $e^{i\pi/3}$. We complement this with a numerical phase diagram for finite $\nu$ (hexagonal $\to$ square $\to$ rectangular in $d=2$, with a certified enclosure of the first transition) and with numerical results on hcp and on mixed pair plus three-body energies.
\end{abstract}
\maketitle

\section{Introduction}
Lattice energy minimisation for \emph{pair} interactions is a classical subject; for completely monotone pair potentials in the plane the hexagonal lattice is optimal \cite{Rankin,Cassels,Ennola,Diananda,Montgomery}, and in higher dimensions universal optimality has been proved in dimensions 8 and 24 \cite{CKMRV}.
Much less is known for genuine many-body energies. Here we study the simplest one, a product-form three-body power law
\begin{equation}\label{eq:T}
T_\nu(\Lat)=\sum_{\substack{x,y\in\Lat\setminus\{0\}\\ x\neq y}}\bigl(|x|\,|y|\,|x-y|\bigr)^{-\nu},
\end{equation}
which converges for $\nu>2d/3$. This sum (under the name three-body zeta function $\zeta^{(3)}_\Lat(\nu,\nu,\nu)$) appears as the radial part of the Axilrod--Teller--Muto energy at $\nu=3$ in \cite{RoblesNavarro}, where it was evaluated for selected lattices and along the Bain path; efficient algorithms are given in \cite{BuchheitBusse,BuchheitRupp}. The reduction of a steep-decay limit to a packing-type problem is standard for pair energies \cite{BST}; three-point max--min problems have been studied on the sphere \cite{Bilyk}, and the criticality of symmetric lattices is automatic \cite{BetCrit}.
We stress that \eqref{eq:T} is a \emph{model} energy: physical three-body dispersion carries an angular factor, and a pure three-body energy at fixed density does not describe a material. The motivation here is mathematical, see Section~\ref{sec:mgf}: in $d=2$, $T_{2s}$ is a modular graph function.

\smallskip\noindent\textbf{What is automatic and what is not.} $T_\nu$ is invariant under isometries, so on the modular surface the hexagonal point $\rho=e^{i\pi/3}$ and the square point $i$ (elliptic fixed points) are always critical points. The content of the results below is \emph{global} minimality and the location of the transitions.

\smallskip\noindent\textbf{Main result.} Let
\[
a_*=\Bigl(\tfrac{1+\sqrt{17}}{8}\Bigr)^{1/4},\qquad P_*=2a_*^3=\sqrt{a_*^2+a_*^{-2}}\approx 1.4317345,\qquad y_\infty=a_*^{-2}=\sqrt{\tfrac{\sqrt{17}-1}{2}}\approx1.2496211,
\]
and let $\Lat_*=a_*\Z\times a_*^{-1}\Z$ (rectangular, aspect ratio $y_\infty$, unit covolume).

\begin{theorem}\label{thm:main}
Let $d=2$. For every $\nu>4/3$ the energy $T_\nu$ attains its minimum on the space of unit-covolume lattices. If $\Lat_\nu$ is any minimiser, then $\Lat_\nu\to\Lat_*$ up to isometries as $\nu\to\infty$, and
$\lim_{\nu\to\infty}\bigl(\min T_\nu\bigr)^{1/\nu}=1/P_*$.
Equivalently (Section~\ref{sec:mgf}), minimisers $\tau_s$ of the modular graph function $C_{s,s,s}$ on the fundamental domain satisfy $\tau_s\to i\,y_\infty$ as $s\to\infty$.
\end{theorem}

Theorem~\ref{thm:main} follows from an exact solution of the following max--min problem.
\begin{definition}
For a lattice $\Lat$ let $\pi(p,q,r)=|p-q|\,|q-r|\,|r-p|$ and
$P(\Lat)=\min\{\pi(p,q,r):p,q,r\in\Lat\ \text{pairwise distinct}\}$ (the minimum exists since $\Lat$ is discrete).
\end{definition}
\begin{theorem}\label{thm:maxmin}
For every lattice $\Lat\subset\R^2$ of covolume $1$, $P(\Lat)\le P_*$, with equality if and only if $\Lat$ is isometric to $\Lat_*$.
\end{theorem}

\section{Proof of Theorem~\ref{thm:maxmin}}
\begin{lemma}[reduced basis]\label{lem:red}
After a rotation and a reflection, every unit-covolume $\Lat\subset\R^2$ has a basis $u=(a,0)$, $v=(p,a^{-1})$ with $a=\lambda_1(\Lat)$, $0\le p\le a/2$ and $|v|\ge a$. Moreover $a^4\le 4/3$.
\end{lemma}
\begin{proof}
Lagrange--Gauss reduction; the bound is Hermite's constant $\gamma_2=2/\sqrt3$.
\end{proof}

\begin{lemma}[upper bound]\label{lem:upper}
With the notation of Lemma~\ref{lem:red},
$P(\Lat)\le G(a):=\min\{2a^3,\sqrt{a^2+a^{-2}}\}$, and $\pi(0,u,v)\le\sqrt{a^2+a^{-2}}$ with equality iff $p=0$.
\end{lemma}
\begin{proof}
The collinear triple $(0,u,2u)$ gives $\pi=a\cdot a\cdot2a=2a^3$. For the triple $(0,u,v)$,
$\pi(0,u,v)^2=a^2\,(p^2+k)\,((a-p)^2+k)$ with $k=a^{-2}$. Put $m=p(a-p)\in[0,a^2/4]$. Since $p^2+(a-p)^2=a^2-2m$,
\[
(p^2+k)((a-p)^2+k)=m^2+k(a^2-2m)+k^2=(ka^2+k^2)+m(m-2k).
\]
Because $a^4\le4/3<8$ we have $m\le a^2/4<2a^{-2}=2k$, hence $m(m-2k)\le0$ with equality iff $m=0$, i.e.\ $p=0$ (note $p\le a/2<a$). Therefore $\pi(0,u,v)^2\le a^2(ka^2+k^2)=a^2+a^{-2}$.
\end{proof}

\begin{lemma}\label{lem:G}
For $0<a\le(4/3)^{1/4}$ we have $G(a)\le P_*$, with equality iff $a=a_*$.
\end{lemma}
\begin{proof}
$a_*$ is the positive root of $4a^8=a^4+1$, equivalently $2a_*^3=\sqrt{a_*^2+a_*^{-2}}$. For $a\le a_*$, $G(a)\le2a^3\le2a_*^3$ with equality iff $a=a_*$. The function $g(a)=a^2+a^{-2}$ is strictly decreasing on $(0,1]$ and increasing on $[1,\infty)$. For $a_*<a\le1$, $G(a)\le\sqrt{g(a)}<\sqrt{g(a_*)}=P_*$. For $1\le a\le(4/3)^{1/4}$, $g(a)\le g((4/3)^{1/4})=\tfrac{2}{\sqrt3}+\tfrac{\sqrt3}{2}=\tfrac{7}{2\sqrt3}$, and $\bigl(\tfrac{7}{2\sqrt3}\bigr)^2=\tfrac{49}{12}<16\bigl(\tfrac{1+\sqrt{17}}{8}\bigr)^3=P_*^4$ (numerically $4.083<4.202$), so again $G(a)<P_*$.
\end{proof}

\begin{lemma}[the extremal lattice]\label{lem:star}
$P(\Lat_*)=P_*$.
\end{lemma}
\begin{proof}
Write $u=(a_*,0)$, $v=(0,a_*^{-1})$; note $a_*<1$. The triples $(0,u,2u)$ and $(0,u,v)$ give $\pi=P_*$, so $P(\Lat_*)\le P_*$.
\emph{Collinear triples.} Three distinct lattice points on a line are $q+j s w,\ q+ksw,\ q+lsw$ with $w$ a unit vector, $s\ge\lambda_1=a_*$ and integers $j<k<l$; then $\pi=s^3(k-j)(l-k)(l-j)\ge2s^3\ge2a_*^3=P_*$.
\emph{Non-collinear triples.} Let the side lengths be $d_1\le d_2\le d_3$ and $\gamma$ the angle between the sides $d_1,d_2$. The triangle has area $\ge1/2$, so $d_1d_2\ge d_1d_2\sin\gamma=2\,\mathrm{area}\ge1$ and $\pi\ge d_3$.
We claim $d_3\ge D:=\sqrt{a_*^2+a_*^{-2}}=P_*$. The three side vectors are nonzero, pairwise linearly independent vectors of $\Lat_*$. A vector $(ma_*,na_*^{-1})$ has norm $<D$ only if $(m,n)\in\{(\pm1,0),(0,\pm1)\}$: if $n\neq0$ then $m^2a_*^2<a_*^2+(1-n^2)a_*^{-2}\le a_*^2$ forces $m=0$ and $n^2<1+a_*^4<2$; if $n=0$ then $m^2<1+a_*^{-4}=1+y_\infty^2<4$. These vectors span only two directions, so at most two of the three pairwise independent side vectors can be shorter than $D$. Hence $d_3\ge D$ and $\pi\ge P_*$.
\end{proof}

\begin{proof}[Proof of Theorem~\ref{thm:maxmin}]
By Lemmas~\ref{lem:red}--\ref{lem:G}, $P(\Lat)\le G(a)\le P_*$. If $P(\Lat)=P_*$, then $G(a)=P_*$, so $a=a_*$; moreover $\pi(0,u,v)\ge P(\Lat)=P_*=\sqrt{a_*^2+a_*^{-2}}$, which by Lemma~\ref{lem:upper} forces $p=0$, i.e.\ $\Lat=\Lat_*$. Conversely $P(\Lat_*)=P_*$ by Lemma~\ref{lem:star}.
\end{proof}

\section{Proof of Theorem~\ref{thm:main}}
Let $\mathcal L_2$ be the space of unit-covolume lattices modulo isometries, with the topology induced by bases.
\begin{lemma}\label{lem:basic}
(a) $T_\nu(\Lat)\ge P(\Lat)^{-\nu}$ for all $\Lat$ and $\nu>4/3$.
(b) $P$ is continuous on $\mathcal L_2$ and $P(\Lat)\le2\lambda_1(\Lat)^3$.
(c) $T_\nu$ is lower semicontinuous on $\mathcal L_2$.
(d) For $\nu\ge\nu_0>4/3$: $T_\nu(\Lat_*)\le C\,P_*^{-\nu}$ with $C=T_{\nu_0}(\Lat_*)\,P_*^{\nu_0}$.
\end{lemma}
\begin{proof}
(a) Translating a minimal triple so that one vertex is $0$ produces a term $P(\Lat)^{-\nu}$ of \eqref{eq:T}.
(b) The bound is the triple $(0,u,2u)$. Continuity: near a given lattice only triples of points in a fixed ball are relevant, and $\pi$ is continuous in the basis.
(c) $T_\nu$ is a sum of nonnegative continuous functions of the basis.
(d) Every term satisfies $\pi\ge P_*$, so $(P_*/\pi)^{\nu}\le(P_*/\pi)^{\nu_0}$.
\end{proof}

\begin{proof}[Proof of Theorem~\ref{thm:main}]
\emph{Existence.} By Lemma~\ref{lem:basic}(a),(b), $T_\nu(\Lat)\ge(2\lambda_1^3)^{-\nu}$, so sublevel sets of $T_\nu$ lie in $\{\lambda_1\ge\varepsilon\}$, which is compact by Mahler's criterion; with (c) a minimiser exists.
\emph{Convergence.} For a minimiser $\Lat_\nu$ and $\nu\ge\nu_0$, Lemma~\ref{lem:basic} gives
$P(\Lat_\nu)^{-\nu}\le T_\nu(\Lat_\nu)\le T_\nu(\Lat_*)\le CP_*^{-\nu}$, hence $P_*\ge P(\Lat_\nu)\ge C^{-1/\nu}P_*\to P_*$.
In particular $\lambda_1(\Lat_\nu)$ is bounded below, so $(\Lat_\nu)$ lies in a compact set. Any limit point $\Lat_\infty$ satisfies $P(\Lat_\infty)=P_*$ by continuity, hence $\Lat_\infty\cong\Lat_*$ by Theorem~\ref{thm:maxmin}. The same chain of inequalities gives $(\min T_\nu)^{1/\nu}\to1/P_*$.
\end{proof}

\section{Relation to modular graph functions}\label{sec:mgf}
For $\tau$ in the upper half plane the two-loop modular graph function is
$C_{a,b,c}(\tau)=\sum_{p_1+p_2+p_3=0,\ p_i\neq0}\prod_{i}\bigl(\tfrac{\operatorname{Im}\tau}{\pi|p_i|^2}\bigr)^{a_i}$, $p_i\in\Z+\tau\Z$ \cite{DGV}.
With $\Lat_\tau=(\Z+\tau\Z)/\sqrt{\operatorname{Im}\tau}$ (unit covolume) one has
\[
C_{s,s,s}(\tau)=\pi^{-3s}\,T_{2s}(\Lat_\tau),
\]
so Theorem~\ref{thm:main} is a statement about the location of the minimum of $C_{s,s,s}$ on the fundamental domain. At $s=1$, Zagier's identity $C_{1,1,1}=E_3+\zeta(3)$ and the minimality of the Epstein zeta function at $\rho$ \cite{Rankin,Cassels,Ennola,Diananda} show that the minimum is at $\rho$ (hexagonal). Theorem~\ref{thm:main} shows that for large $s$ it is near $i y_\infty$. We are not aware of earlier results on the minima of $C_{s,s,s}$ for $s>1$; see the caveat in Section~\ref{sec:limits}.

\section{Finite \texorpdfstring{$\nu$}{nu}: the minimum of \texorpdfstring{$C_{s,s,s}$}{C_sss} for \texorpdfstring{$s=2,3,4$}{s=2,3,4}}\label{sec:finite}
\begin{theorem}[computer-assisted]\label{thm:sss}
For $s\in\{2,3,4\}$ the modular graph function $C_{s,s,s}$ attains its minimum over the fundamental domain only at $\tau=i$. Equivalently, for $\nu\in\{4,6,8\}$ the square lattice is the unique minimiser of $T_\nu$ among unit-covolume lattices in $\R^2$. Together with the case $s=1$ (minimum at $\rho$, Section~\ref{sec:mgf}) and Theorem~\ref{thm:main}, the minimiser of $C_{s,s,s}$ is $\rho$ for $s=1$, $i$ for $s=2,3,4$, and tends to $i\,y_\infty$ as $s\to\infty$.
\end{theorem}
\noindent\emph{A universal majorant.} Write $\tau=x+iy$, $\ell_v=|m+n\tau|^2/y$ for $v=(m,n)$, so that $T_\nu=\sum\prod_{\text{sides}}\ell^{-s}$. Around $(x_0,y_0)$ put $A=a/y_0$, $B=b/y_0$. A short computation gives
\[
\frac{\ell_v(x_0+a,y_0+b)}{\ell_v(x_0,y_0)}=1+\frac{(2\alpha^2-1)B+2\alpha\beta A+\alpha^2(A^2+B^2)}{1+B},
\]
with $\alpha^2+\beta^2=1$ depending on $v$. Hence, coefficientwise, $\ell/\ell_0-1\ll\hat\varepsilon:=(A+B+A^2+B^2)/(1-B)$ and every triangle term $t$ satisfies $t/t_0\ll(1-\hat\varepsilon)^{-3s}$, \emph{uniformly in the lattice vectors}. Consequently all Taylor coefficients of $T_\nu$ at $\tau_0$ are bounded by $T_\nu(\tau_0)$ times the coefficients $\hat G_{ij}$ of $(1-\hat\varepsilon)^{-3s}$; this yields explicit remainder bounds and explicit bounds for the contribution of the terms omitted by a finite truncation.

\noindent\emph{Proof outline.} (1) For $y>Y_0$ the collinear triples along the vector $1$ give $T_\nu(\tau)\ge y^{3s}\sum_{j\ne k}(|j||k||j-k|)^{-2s}>T_\nu(i)$. (2) Local lemma: the symmetries $x\mapsto-x$ and $\tau\mapsto1/\bar\tau$ fix $i$ and force $c_{10}=c_{01}=c_{11}=c_{30}=c_{12}=0$ in the Taylor expansion at $i$; with rigorous enclosures of $c_{20},c_{02},c_{21},c_{03}$ and the majorant remainder, $T(i+\delta)-T(i)>0$ for $0<\max(|A|,|B|)\le r_{\rm loc}$. (3) The rest of $[0,\tfrac12]\times[0.8,Y_0]$ is covered by boxes on which the second-order Taylor polynomial at the centre (with rigorous coefficient enclosures) minus the majorant remainder exceeds an upper bound for $T(i)$. Coefficients are computed from cube-truncated sums by direct (non-FFT) convolution with an explicit tail bound and a rounding margin.

\noindent\emph{Data of the three runs} (scripts \texttt{theorem\_mgf/prove\_square\_min.py}): 
\begin{center}\small
\begin{tabular}{cccccc}\toprule
$\nu$ & $T_\nu(i)$ (enclosure) & $c_{20}\ge$ & $c_{02}\ge$ & $r_{\rm loc}$ & boxes\\\midrule
4 & $[7.5826697721,\,7.5826700248]$ & 0.43688 & 9.75784 & 0.00309 & 531\\
5 & $[4.8380268461,\,4.8380268473]$ & 2.58218 & 6.30394 & 0.00701 & 303\\
6 & $[3.2499743121,\,3.2499743128]$ & 4.27194 & 3.57092 & 0.00755 & 307\\
7 & $[2.2334461746,\,2.2334461750]$ & 5.34452 & 1.66169 & 0.00495 & 323\\
8 & $[1.5524012124,\,1.5524012128]$ & 5.83081 & 0.46080 & 0.00260 & 385\\\bottomrule
\end{tabular}\end{center}
The runs for $\nu=5,7$ show that the square lattice is also the unique minimiser of $T_5$ and $T_7$. The same approach failed for the hexagonal lattice at $\nu=3$: close to $\nu=d$ the truncation tail decays only like $R^{-4}$, and at $R=40$ the error exceeds the energy gap next to $\rho$. The small coefficient $c_{02}$ at $\nu=8$ reflects the nearby loss of stability of the square lattice at $\nu_2\approx8.606$.
The status of the floating-point parts is the same as in Remark after Theorem~\ref{thm:3d}: sums are sequential double-precision sums with explicit margins, not interval arithmetic.

\section{Numerical phase diagram for finite \texorpdfstring{$\nu$}{nu}}\label{sec:numerics}
The following are numerical observations, not theorems. $T_\nu$ was evaluated by cube-truncated direct summation with an FFT convolution and Richardson extrapolation in the cut-off (independent of \cite{BuchheitRupp}); agreement with the library of \cite{BuchheitRupp} is better than $2\cdot10^{-10}$ relative on 55 test cases, and the values at $\nu=3$ reproduce those of \cite{RoblesNavarro} after conversion to unit covolume.
\begin{itemize}
\item $2<\nu<\nu^*$: the hexagonal lattice is the global minimiser. At $\nu=2$ this is exact by Section~\ref{sec:mgf}.
\item $\nu^*=3.9183649026\ldots$: first-order transition to the square lattice; both are local minima for $3.806<\nu<4.278$. \emph{Certified:} $T_\nu(\Z^2)-T_\nu(\mathrm{hex})$ changes sign in $(3.918364,3.918366)$; this uses a direct (non-FFT) evaluation of the cube-truncated sum at $R=100$, an explicit tail bound, and a rounding margin (script \texttt{certify\_nustar.py}).
\item $\nu_2\approx8.606$: the square lattice becomes unstable towards rectangular deformations; the optimal aspect ratio then grows continuously (1.056 at $\nu=9$, 1.100 at $10$, 1.141 at $12$, 1.170 at $15$, 1.194 at $20$, 1.215 at $30$), consistent with the limit $y_\infty=1.2496$ of Theorem~\ref{thm:main}.
\end{itemize}
\begin{conjecture}
For $d=2$ the minimiser of $T_\nu$ is hexagonal for $4/3<\nu<\nu^*$, square for $\nu^*<\nu\le\nu_2$ and rectangular with aspect ratio increasing from $1$ to $y_\infty$ for $\nu>\nu_2$.
\end{conjecture}

\section{Three dimensions}\label{sec:3d}
For unit-covolume $\Lat\subset\R^3$ the energy $T_\nu$ converges for $\nu>2$, and the argument of Theorem~\ref{thm:main} applies verbatim once the max--min problem is solved.
For the body-centred cubic lattice $P(\mathrm{BCC})=3/2$, attained by the isosceles triangles with two nearest-neighbour sides and one second-neighbour side, while $P(\mathrm{FCC})=P(\mathrm{SC})=\sqrt2$ (FCC: equilateral nearest-neighbour triangles; SC: right isosceles triangles).
\begin{theorem}[computer-assisted]\label{thm:3d}
For every lattice $\Lat\subset\R^3$ of covolume $1$, $P(\Lat)\le 3/2$, with equality if and only if $\Lat$ is isometric to the BCC lattice of covolume $1$. Consequently, for $\nu>2$ the energy $T_\nu$ has minimisers, and every family of minimisers converges to BCC as $\nu\to\infty$, with $(\min T_\nu)^{1/\nu}\to 2/3$.
\end{theorem}
\noindent\emph{Scale-invariant formulation.} Write $G=B^{\mathsf T}B$ for the Gram matrix of a basis $B$ and $Q(G)=P(\Lat)/\sqrt{\det G}$. By scaling we may assume $G_{11}=\lambda_1^2=1$ and use the parameters $q=(g_{22},g_{33},g_{12},g_{13},g_{23})$.

\begin{lemma}[compact search region]\label{lem:region3}
Every lattice has, up to isometry and scaling, a Gram matrix in the region $\mathcal R$ defined by
$1\le g_{22}\le g_{33}$, $0\le g_{12},g_{13}\le \tfrac12$, $|g_{23}|\le g_{22}/2$, $1+g_{22}+2(\varepsilon_1\varepsilon_2g_{12}+\varepsilon_1g_{13}+\varepsilon_2g_{23})\ge0$ for $\varepsilon_i=\pm1$.
On $\mathcal R$ one has $\det G\ge g_{22}g_{33}/4$, hence $Q(G)<3/2$ whenever $g_{22}g_{33}>64/9$.
\end{lemma}
\begin{proof}
Take a Minkowski-reduced basis and flip the signs of $b_2,b_3$ to make $g_{12},g_{13}\ge0$; the listed inequalities are part of Minkowski's conditions. For the determinant: the triangle $(0,b_1,b_2)$ is acute, so the covering radius $\rho$ of the plane lattice $L_2=\Z b_1+\Z b_2$ is its circumradius, $\rho^2=g_{22}|b_2-b_1|^2/(4\det_2)$ with $\det_2=g_{22}-g_{12}^2\ge\tfrac34g_{22}$, whence $\rho^2\le(1+g_{22})/3$. Minkowski reduction implies that the projection of $b_3$ to $\operatorname{span}(b_1,b_2)$ lies in the Voronoi cell of $L_2$, so the height $h$ of $b_3$ above that plane satisfies $h^2\ge g_{33}-\rho^2\ge g_{33}/3$. Thus $\det G=\det_2h^2\ge g_{22}g_{33}/4$. Finally the collinear triple $(0,b_1,2b_1)$ gives $Q\le2/\sqrt{\det G}\le4/\sqrt{g_{22}g_{33}}$.
\end{proof}
In particular it suffices to consider the compact box $g_{22}\in[1,8/3]$, $g_{33}\in[1,64/9]$, $g_{12},g_{13}\in[0,1/2]$, $|g_{23}|\le4/3$.
Within $\mathcal R$ BCC has exactly one representative, $q_0=(1,1,\tfrac13,\tfrac13,-\tfrac13)$ (checked by enumerating unimodular changes of basis).

\begin{lemma}[local lemma, interval arithmetic]\label{lem:local3}
Let $f_t(q)=\sqrt{\ell_{m}\ell_{n}\ell_{m-n}/\det G}$ for a triangle $t=(0,m,n)$, where $\ell_v=v^{\mathsf T}Gv$. At $q_0$ exactly $36$ triangles (up to order) are active, $f_t(q_0)=3/2$, and their gradients have $6$ distinct values whose convex hull contains the origin at distance $c=0.13887\ldots$ from its boundary. On the box $q_0+[-r,r]^5$ with $r=0.007$, interval arithmetic gives $\|\nabla^2f_t\|_F\le M=14.92$ for all active $t$. Hence $\min_tf_t(q_0+\delta)\le\tfrac32-c|\delta|+\tfrac M2|\delta|^2<\tfrac32$ for $0<|\delta|\le\sqrt5\,r=0.01565<2c/M=0.01862$.
\end{lemma}

\noindent\emph{Branch and bound.} On a box $B$ of parameters, each $\ell_v$ is linear in $q$ and $\det G$ is a cubic polynomial; we bound $Q\le\min_{t\in S}\bigl(\prod\max_B\ell\bigr)^{1/2}/(\min_B\det G)^{1/2}$ for a set $S$ of candidate triangles (all $(0,m,n)$ with $m,n\in\{-1,0,1\}^3$, and the collinear $(0,m,2m)$), with a relative safety margin $10^{-9}$ for rounding. Starting from the compact box above, boxes are discarded when they violate a region constraint, lie inside $q_0+[-r,r]^5$, or have upper bound $<3/2$, and are bisected otherwise. The run terminates after $9.93\cdot10^6$ boxes: \emph{every} box is certified.
Together with Lemmas~\ref{lem:region3} and~\ref{lem:local3} this gives $Q\le3/2$ on $\mathcal R$ with equality only at $q_0$, i.e.\ Theorem~\ref{thm:3d}. The convergence statement follows as in the proof of Theorem~\ref{thm:main}.

\begin{remark}[status of the computer-assisted parts]
The local lemma uses rigorous interval arithmetic (\texttt{mpmath.iv}). The branch and bound uses IEEE double precision with explicit safety margins ($10^{-9}$ relative, $10^{-12}$ absolute on the determinant) that exceed the rounding errors of the few arithmetic operations per bound by several orders of magnitude; a fully interval-arithmetic re-run is straightforward and is planned. As sanity checks, the box bounds were compared with brute-force values of $Q$ at 1200 random points (never violated), and the prover correctly fails when asked to certify $Q<1.49$.
\end{remark}

\subsection*{Finite $\nu$ in three dimensions (numerical)}
BCC is the best lattice found by a global search over all Bravais lattices for every tested $\nu\in[3.2,12]$. FCC is a local minimum only for small $\nu$: its curvature along the Bain path changes sign at $\nu_c=3.7521\pm0.0001$ (spread over finite-difference steps and cut-offs). Among non-Bravais structures, hcp (with optimised $c/a$) lies slightly above FCC for all tested $\nu$; for $\nu\to\infty$ this is also clear from $P(\mathrm{hcp})=P(\mathrm{FCC})=\sqrt2<3/2$.
For mixed energies $E_\lambda=Z_\nu+\lambda T_\nu$ (pair Riesz energy $Z_\nu=\sum_{x\ne0}|x|^{-\nu}$ plus three-body term, unit density) the optimum among FCC, BCC, hcp$(c/a)$ and the Bain family jumps directly from FCC to BCC at a critical ratio $\lambda_c(\nu)$ (e.g.\ $\lambda_c\approx0.05$, $0.07$, $0.11$, $0.32$ for $\nu=3.5,4.5,6,9$); hcp is never optimal.

\section{Limitations}\label{sec:limits}
(i) $T_\nu$ is a model energy; physically relevant three-body dispersion has an angular factor and is always accompanied by pair terms. (ii) The theorems concern Bravais lattices; hcp was only compared numerically. (iii) Apart from Theorems~\ref{thm:main}, \ref{thm:maxmin}, \ref{thm:3d}, \ref{thm:sss} and the certified enclosure of $\nu^*$, all statements are numerical; in particular the finite-$\nu$ phase diagrams are conjectural. (iv) The computer-assisted parts use double precision with explicit safety margins rather than interval arithmetic throughout. (v) Our literature search for prior results on minima of $C_{s,s,s}$ and on minimisers of $T_\nu$ was not exhaustive.

\begin{thebibliography}{99}
\bibitem{Rankin} R.\,A.~Rankin, A minimum problem for the Epstein zeta-function, Proc. Glasgow Math. Assoc. 1 (1953).
\bibitem{Cassels} J.\,W.\,S.~Cassels, On a problem of Rankin about the Epstein zeta-function, Proc. Glasgow Math. Assoc. 4 (1959).
\bibitem{Ennola} V.~Ennola, A lemma about the Epstein zeta-function, Proc. Glasgow Math. Assoc. 6 (1964).
\bibitem{Diananda} P.\,H.~Diananda, Notes on two lemmas concerning the Epstein zeta-function, Proc. Glasgow Math. Assoc. 6 (1964).
\bibitem{Montgomery} H.\,L.~Montgomery, Minimal theta functions, Glasgow Math. J. 30 (1988).
\bibitem{CKMRV} H.~Cohn, A.~Kumar, S.\,D.~Miller, D.~Radchenko, M.~Viazovska, Universal optimality of the $E_8$ and Leech lattices and interpolation formulas, Ann. of Math. 196 (2022).
\bibitem{RoblesNavarro} A.~Robles-Navarro et al., arXiv:2504.07338, J. Chem. Phys. 163 (2025) 094104.
\bibitem{BuchheitBusse} A.\,A.~Buchheit, J.\,T.~Busse, arXiv:2504.11989.
\bibitem{BuchheitRupp} A.\,A.~Buchheit, A.~Rupp, arXiv:2609.18918.
\bibitem{DGV} E.~D'Hoker, M.\,B.~Green, P.~Vanhove, On the modular structure of the genus-one Type II superstring low energy expansion, JHEP 08 (2015) 041, arXiv:1502.06698.
\bibitem{BST} L.~B\'etermin, L.~\v{S}amaj, I.~Trav\v{e}nec, arXiv:2107.14020.
\bibitem{Bilyk} D.~Bilyk et al., arXiv:2303.12283.
\bibitem{BetCrit} L.~B\'etermin, arXiv:1611.07798.
\end{thebibliography}
\end{document}
